\documentclass[12pt,a4paper]{article}
\usepackage[utf8]{inputenc}
\usepackage[T1]{fontenc}
\usepackage{lmodern}
\usepackage{amsmath, amssymb, amsthm, amsfonts, mathrsfs, graphicx, physics, enumitem, geometry, esint,twemojis}
\geometry{margin=1in}
\usepackage[dvipsnames]{xcolor}
\definecolor{EvanPink}{rgb}{0.85, 0.13, 0.55}
\usepackage[colorlinks=true,
            linkcolor=OliveGreen,
            urlcolor=EvanPink,
            citecolor=OliveGreen]{hyperref}
\usepackage{cleveref}
\usepackage{etoolbox}

\newenvironment{solution}{%
  \par\noindent\textit{Solution.}\ }{\qed}

\newtheoremstyle{problemstyle}
  {1em}   
  {1em}   
  {}      
  {}      
  {\bfseries} 
  {.}     
  {0.5em} 
  {}      
\theoremstyle{problemstyle}
\newtheorem{problem}{Problem}[section]

\apptocmd{\problem}{%
  \addcontentsline{toc}{subsection}{Problem~\thesection.\arabic{problem}}%
}{}{}

\crefname{problem}{Problem}{Problems}
\Crefname{problem}{Problem}{Problems}

\setcounter{tocdepth}{2}

\title{\textbf{Chapter 3 : Signed Measures and Differentiation}}
\author{Arkaraj Mukherjee}

\def\bN{\mathbb{N}}
\def\bR{\mathbb{R}}
\def\bZ{\mathbb{Z}}
\def\bQ{\mathbb{Q}}
\def\bC{\mathbb{C}}
\def\mM{\mathcal{M}}
\def\mR{\mathcal{R}}
\def\mN{\mathcal{N}}
\def\mE{\mathcal{E}}
\def\mA{\mathcal{A}}
\def\oR{\overline{\mathbb{R}}}

\begin{document}
\maketitle
\newpage
\tableofcontents
\newpage
\setcounter{section}{2}
\stepcounter{section}
\section*{Chapter \thesection\ : Signed Measures and Differentiation}
\addcontentsline{toc}{section}{Chapter \thesection\ : Signed Measures and Differentiation}

\begin{problem}\label{prob:3.1}
Prove Proposition 3.1.
\end{problem}
\begin{solution}
Wlog suppose that $\nu$ omits the value $-\infty$, then if $E_n$ is a sequence increasing sets we have that,
\[
	\nu \qty(\bigcup_{n=1}^\infty E_n)=\nu(E_1)+ \sum_{n\geq 2}\nu(E_n\setminus E_{n-1})=\lim_{m\to\infty} \nu(E_1)+ \sum_{n\leq m}\nu(E_n\setminus E_{n-1})=\lim_{m\to\infty}\nu(E_m)
\]
And if it was decreasing with $\nu(E_1)$ finite then, $E_1\setminus E_n$ is increasing and we have
\[
	\nu(E_1)-\nu \qty(\bigcap E_n)=\nu \qty(E_1\setminus\bigcap E_n)=\nu \qty(\bigcup(E_1\setminus E_n))=\lim_{m\to\infty}\nu(E_1\setminus E_m)=\nu(E_1)- \lim_{m\to\infty}\nu(E_m)
\]
after which we can cancel out the finite term of $\nu(E_1)$.
\end{solution}
\begin{problem}\label{prob:3.2}
If $\nu$ is a signed measure, $E$ is $\nu$-null iff $|\nu|(E)=0$. Also, if $\nu$ and $\mu$ are signed measures, $\nu\perp\mu$ iff $|\nu|\perp\mu$ iff $\nu^{+}\perp\mu$ and $\nu^{-}\perp\mu$.
\end{problem}
\begin{solution}
If $|\nu|(E)=0$ then $\nu^+(E)=\nu^-(E)=0$ and as these are positive measures, by monotonicity see that $\nu^+(F)=\nu^-(F)=0$ for any measurable $F\subseteq E$ making $E$ $\nu$-null. 
Now, if $E$ was $\nu$-null then $|\nu|(E)>0$ would imply that $\nu^+(E)$ or $\nu^-(E)$ is strictly greater than $0$ and if we assume the first to be as described then $\nu(E\cap P)=\nu^+(E\cap P)>0$ which contradicts nullity of $E$ w.r.t $\nu$ as $E\cap P \subseteq E$. 
The other case leads to a contradiction similarly, thus we must have $|\nu|(E)=0$. 

As a measurable set $E$ is $|\nu|$-null iff $|\nu|(E)=0$ stemming from positivity we see that $\nu\perp\mu$ iff $|\nu|\perp\mu$ by the previous claim. 
We also see that $|\nu|\perp\mu$ iff $\nu^+,\nu^-\perp\mu$ as $|\nu|(E)=0$ iff $\nu^-(E),\nu^+(E)=0$, which proves the equivalence.
\end{solution}
\begin{problem}\label{prob:3.3}
Let $\nu$ be a signed measure on $(X, \mM)$.
\begin{enumerate}[label=\alph*.]
    \item $L^{1}(\nu)=L^{1}(|\nu|)$.
    \item If $f\in L^{1}(\nu)$, $|\int f~d\nu|\le\int|f|d|\nu|$.
    \item If $E\in\mM$, $|\nu|(E) = \sup \{ |\int_E f\, d\nu| : |f| \le 1 \}$.
\end{enumerate}
\end{problem}
\begin{solution}
\begin{enumerate}[label=\alph*.]
\item As $L^1(\nu):=L^1(\nu^+)\cap L^1(\nu^-)$ we find that, $f\in L^1(\nu)$ implies $ \int_{}^{}|f|,\mathrm{d}\nu^+, \int_{}^{}|f|\,\mathrm{d}\nu^-<\infty$ and thus, $ \int_{}^{}|f|\,\mathrm{d}|\nu|= \int_{}^{}|f|\,\mathrm{d}\nu^++ \int_{}^{}|f|\,\mathrm{d}\nu^-<\infty$ where we have used the following lemma: If $\mu_1,\mu_2$ are positive measures on the same sigma algebra then so is their sum and if $f\in L^+$ then $ \int_{}^{}f\,\mathrm{d}(\mu_1+\mu_2)= \int_{}^{}f\,\mathrm{d}\mu_1+ \int_{}^{}f\,\mathrm{d}\mu_2$. 
For a proof, one can look at simple functions first where its trivial and then routinely extend to $L^+$ by MCT and approximations using simple functions. 
Now, if we had $f\in L^1(|\nu|)$ then as $ \int_{}^{}f\,\mathrm{d}|\nu|=\int_{}^{}|f|\,\mathrm{d}\nu^++ \int_{}^{}|f|\,\mathrm{d}\nu^-<\infty$, both of these terms are individually finite which puts $f$ in $L^1(\nu^+)$ and $L^1(\nu^-)$ proving the claim.
\item \[
    \qty|\int fd\nu|=\qty|\int fd\nu^+-\int fd\nu^-|
\]
\[\leq\qty|\int fd\nu^+|+\qty|\int fd\nu^-|\leq\int|f|d\nu^++\int|f|d\nu^-=\int|f|d|\nu|
\]
 
\item Consider the hahn decomposition $P\cup N$ is usual notation leading upto the jordan decomposition, then letting $f=\chi_P-\chi_N$ we see that $ \int_{}^{}f\,\mathrm{d}\nu=\nu(P\cap E)-\nu(N\cap E)=\nu^+(E)+\nu^-(E)=|\nu|(E)$ so we have one direction of the inequality ($\leq$) and for the other direction, by part (b) we can say that $|f|\leq 1$ implies $|\int_{}^{}f\,\mathrm{d}\nu|\leq \int_{}^{}1\,\mathrm{d}|\nu|=|\nu|(E).$ so we are done.
\end{enumerate}
\end{solution}
\begin{problem}\label{prob:3.4}
If $\nu$ is a signed measure and $\lambda$, $\mu$ are positive measures such that $\nu=\lambda-\mu$, then $\lambda\ge\nu^{+}$ and $\mu\ge\nu^{-}$.
\end{problem}
\begin{solution}
Consider the hahn decomposition $P\cup N$ in usual notation, leading upto the jordan decomposition. 
Then we have $\nu^+(E)=\nu(E\cap P)=\lambda(E\cap P)-\mu(E\cap P)\leq\lambda(E\cap P)\leq\lambda(E)$ and the other inequality follows similarly. 
\end{solution}
\begin{problem}\label{prob:3.5}
If $\nu_{1}$, $\nu_{2}$ are signed measures that both omit the value $+\infty$ or $-\infty$, then $|\nu_{1}+\nu_{2}|\le|\nu_{1}|+|\nu_{2}|$. (Use Exercise 4.)
\end{problem}
\begin{solution}
It can be easily shown that in this case $\nu_1+\nu_2$ is also a signed measure. 
As all $\nu_1,\nu_2$ are finite, we can rearrange to get $\nu_1+\nu_2=\nu_1^+-\nu_1^-+\nu_2^+-\nu_2^-=(\nu_1^++\nu_2^+)-(\nu_1^-+\nu_2^-)$ which, by the previous problem implies $(\nu_1+\nu_2)^+\leq\nu_1^++\nu_2^+$ and $(\nu_1+\nu_2)^-\leq\nu_1^-+\nu_2^-$ implying $|\nu_1+\nu_2|\leq(\nu_1^++\nu_2^+)+(\nu_1^-+\nu_2^-)=(\nu_1^++\nu_1^-)+(\nu_2^++\nu_2^-)\leq|\nu_1|+|\nu_2|$.
\end{solution}
\begin{problem}\label{prob:3.6}
Suppose $\nu(E)=\int f~d\mu$ where $\mu$ is a positive measure and $f$ is an extended $\mu$-integrable function. Describe the Hahn decompositions of $\nu$ and the positive, negative, and total variations of $\nu$ in terms of $f$ and $\mu$.
\end{problem}
\begin{solution}
	Let $P= \qty{f>0}$ and $N= \qty{f\leq 0}$, then $P\cup N$ is clearly a hahn decomposition and we see that $\nu^+(E)= \int_{E\cap P}^{}f\,\mathrm{d}\mu= \int_{E}^{}f^+\,\mathrm{d}\mu$ and similarly $\nu^-(E)= \int_{E}^{}f^-\,\mathrm{d}\mu$ and thus, $|\nu|(E)= \int_{E}^{}|f|\,\mathrm{d}\mu$.
\end{solution}
\begin{problem}\label{prob:3.7}
Suppose that $\nu$ is a signed measure on $(X, \mM)$ and $E\in\mM$.
\begin{enumerate}[label=\alph*.]
    \item $\nu^{+}(E)=\sup\{\nu(F):F\in\mM,F\subset E\}$ and $\nu^{-}(E)=-\inf\{\nu(F):F\in\mM, F\subset E\}$.
    \item $|\nu|(E)=\sup\{\sum_{1}^{n}|\nu(E_{j})|:n\in\bN,E_{1},...,E_{n} \text{ are disjoint, and } \bigcup_{1}^{n}E_{j}=E\}$.
\end{enumerate}
\end{problem}
\begin{solution}
\begin{enumerate}[label=\alph*.]
	\item $F \subseteq E$ implies that $\nu(F)\leq\nu^+(F)\leq\nu^+(E)$ and $\nu(E\cap P)=\nu^+(E)$ with $E\cap P \subseteq E$ proving one part, the other follows similarly.
	\item Considering a hahn decomposition $P\cup N$ leading to a jordan decomposition, we see that $|\nu|(E)=\nu^+(E)+\nu^-(E)=|\nu(E\cap P)|+|\nu(E\cap N)|$ as $\nu^+(E)=\nu(E\cap P)\geq 0$ and $\nu^-(E)=-\nu(E\cap N)\geq 0$ proving one direction of the equality. 
	For the other, it can be easily proved that $|\nu(E)|\leq|\nu|(E)$ and thus, $\bigcup_1^nE_j=E$ (disjoint union) implies $\sum_1^n|\nu(E_j)|\leq\sum_1^n|\nu|(E_j)=|\nu|(E)$ so we are done.
\end{enumerate}
\end{solution}
\begin{problem}\label{prob:3.8}
$\nu\ll\mu$ iff $|\nu|\ll\mu$ iff $\nu^{+}\ll\mu$ and $\nu^{-}\ll\mu$.
\end{problem}
\begin{solution}
If $\nu\ll\mu$ then $\mu(E)=0\implies\nu(E)=0$ and as $\mu$ is taken positive we must have that $\mu(E\cap P),\mu(E\cap N)=0$ where $P\cup N$ is the hahn decomposition for $\nu$, this implies that $\nu(E\cap P),\nu(E\cap N)=0$ so $\nu^\pm(E)=0$ and we have that $|\nu|(E)=0$ as well proving $|\nu|\ll\mu$. 
As $|\nu|(E)\geq|\nu(E)|$ we have that $|\nu|\ll\mu$ implies $\nu\ll\mu$ and hence $\nu\ll\mu$ iff $|\nu|\ll\mu$. 
If $\nu^+,\nu^-\ll\mu$ then $\mu(E)=0$ clearly implies $|\nu|(E)=0$ and if we have $|\nu|\ll\mu$ then $\mu(E)=0$ implies $\nu^+(E)+\nu^-(E)=0$ and thus $\nu^+(E)=\nu^-(E)=0$ proving $\nu^+,\nu^-\ll\mu$ and thus all these conditions are equivalent.
\end{solution}
\begin{problem}\label{prob:3.9}
Suppose $\{\nu_{j}\}$ is a sequence of positive measures. If $\nu_{j}\perp\mu$ for all $j$, then $\sum_{1}^{\infty}\nu_{j}\perp\mu$; and if $\nu_{j}\ll\mu$ for all $j$, then $\sum_{1}^{\infty}\nu_{j}\ll\mu$. (This is also true for complex measures when $\sum_1^\infty\nu_j$ is a valid complex measure too)
\end{problem}
\begin{solution}
Say $E_j$ is $\nu_j$-null and $E_j^c$ is $\mu$-null for all $j$, then we see that, let $E=\bigcap E_j$, then we clearly see that $E$ is $\sum_1^\infty\nu_j$-null and we claim that, $E^c=\bigcup E_j^c$ is $\mu$-null. 
For any measurable $F\subseteq E^c$, we can write $F=\bigcup_{j=1}^\infty(F\cap(E_j^c\setminus\cup_{i\neq j}^\infty E_i^c))$ is a disjoint union and 
\[
    \mu(F)= \sum_{j=1}^{\infty}\mu(F\cap(E_j^c\setminus\cup_{i\neq j}^\infty E_i^c))=0
\]
as $(F\cap(E_j^c\setminus\cup_{i\neq j}^\infty E_i^c)) \subseteq E_j^c$. 
The absolute continuity result is trivial. 
\end{solution}
\begin{problem}\label{prob:3.10}
Theorem 3.5 may fail when $\nu$ is not finite. (Consider $d\nu(x)=dx/x$ and $d\mu(x)=dx$ on $(0, 1)$, or $\nu = \text{counting measure}$ and $\mu(E)=\sum_{n\in E}2^{-n}$ on $\bN$.)
\end{problem}
\begin{solution}
We choose the second counterexample. 
Clearly we have, $\nu\ll\mu$ but for any $\delta>0$ we see that there exists a $E_\delta$ with $\mu(E_\delta)<\delta$ but $\nu(E_\delta)\geq 1$ which just reads ``$E$ is nonempty'' and as the sum $\sum_{n\in\bN}2^{-n}$ converges absolutely we can find such a $E_\delta$ always.
\end{solution}
\begin{problem}\label{prob:3.11}
Let $\mu$ be a positive measure. A collection of functions $\{f_{\alpha}\}_{\alpha\in A}\subset L^{1}(\mu)$ is called uniformly integrable if for every $\epsilon>0$ there exists $\delta>0$ such that $|\int_{E}f_{\alpha}d\mu|<\epsilon$ for all $\alpha\in A$ whenever $\mu(E)<\delta$.
\begin{enumerate}[label=\alph*.]
    \item Any finite subset of $L^{1}(\mu)$ is uniformly integrable.
    \item If $\{f_{n}\}$ is a sequence in $L^{1}(\mu)$ that converges in the $L^{1}$ metric to $f\in L^{1}(\mu)$, then $\{f_{n}\}$ is uniformly integrable.
\end{enumerate}
\end{problem}
\begin{solution}
\begin{enumerate}[label=\alph*.]
	\item Consider $A=\qty{1,\ldots,n}$; for any $\epsilon>0$, as $f_jd\mu\ll\mu$ where $f_jd\mu$ are finite meausres, we can find $\delta_j>0$ such that $|\int_{E}f_jd\mu|<\epsilon$ for all $j\in A$ whenever $\mu(E)<\delta_j$ so taking $\delta=\min_j\delta_j$ suffices.
	\item As $f_n$ converges in $L^1$, it is cauchy in this metric too so for any $\epsilon>0$ we can find $N$ such that $m,n\geq N\implies|\int_Ef_m-\int_Ef_n|\leq\int|f_m-f_n|<\epsilon$ and by the previous problem, as $\{f_1,\ldots,f_N\}$ are uniformly integrable, we can find $\delta>0$ such that, $|\int_E f_j|<\epsilon<2\epsilon$ for $1\leq j\leq N$ whenever $\mu(E)<\delta$ and for $n>N,$
		\[
		    \qty|\int_Ef_nd\mu|\leq\qty|\int_Ef_Nd\mu|+\qty|\int_Ef_n-f_Nd\mu|<\epsilon+\epsilon=2\epsilon
		\]
	so we are done.
\end{enumerate}
\end{solution}
\begin{problem}\label{prob:3.12}
For $j=1,2,$ let $\mu_{j}$, $\nu_{j}$ be $\sigma$-finite measures on $(X_{j},\mM_{j})$ such that $\nu_{j}\ll\mu_{j}$. Then $\nu_{1}\times\nu_{2}\ll\mu_{1}\times\mu_{2}$ and
\[\frac{d(\nu_{1}\times\nu_{2})}{d(\mu_{1}\times\mu_{2})}(x_{1},x_{2})=\frac{d\nu_{1}}{d\mu_{1}}(x_{1})\frac{d\nu_{2}}{d\mu_{2}}(x_{2}).\]
\end{problem}
\begin{solution}
Suppose $\mu_1\times\mu_2(E)=0$ then we see that $0=\int\mu_1(E_x)d\mu_2(x)$ implying $\mu_1(E_x)=0$ (this in turn implies $\nu_1(E_x)=0$ here) for $\mu_2$-a.e. $x$ so suppose its true for all $x$ outside $N$ with $\mu_2(N)=0$ which implies $\nu_2(N)=0$ so we have that $\nu_1\times\nu_2(E)=\int\nu_1(E_x)d\nu_2(x)=0$ as $\nu_1(E_x)=0$ for $\nu_2$-a.e. $x$ thus proving $\nu_1\times\nu_2\ll\mu_1\times\mu_2.$ 
CLearly these product measures are sigma finite too, so we can find a radon nikodym derivative $f$ i.e. $d(\nu_1\times\nu_2)=fd(\mu_1\times\mu_2)$ and also $f_1,f_2$ such that $f_jd\mu_j=\nu_j$ for $j=1,2$ and by fubini's we see that,
\[
	\int\chi_E(x,y)f_1(x)f_2(y)d(\mu_1\times\mu_2(x,y))=\int\int f_1(x)\chi_{E_x}(y)f_2(y)d\mu_2(y)d\mu_1(x)
\]
\[
    =\int\nu_2(E_x)f_1(x)d\mu_1(x)=\int\nu_2(E_x)d\nu_1(x)=\nu_1\times\nu_2(E)=\int\chi_Ef(x,y)d(\mu_1\times\mu_2(x,y))
\]
which proves that $f(x,y)=f_1(x)f_2(y)$ for $\mu_1\times\mu_2$ a.e. $(x,y)$ so we have the equality treating them as equivalence classes.
\end{solution}
\begin{problem}\label{prob:3.13}
Let $X=[0,1]$, $\mM=\mathcal{B}_{[0,1]}$, $m=$ Lebesgue measure, and $\mu=$ counting measure on $\mM$.
\begin{enumerate}[label=\alph*.]
    \item $m\ll\mu$ but $dm\ne f~d\mu$ for any $f$.
    \item $\mu$ has no Lebesgue decomposition with respect to $m$.
\end{enumerate}
\end{problem}
\begin{solution}
\begin{enumerate}[label=\alph*.]
	\item As $\mu(E)=0\implies E=\emptyset$, $m\ll\mu$ holds. 
		If we had $dm=fd\mu$ then we would have $m( \qty{x})=f(x)=0$ for all $x$ which is absurd as it implies $dm=0$. 
    \item If we had $\mu=\rho+\lambda$ where $\lambda\perp m$ and $\rho\ll m$, notice how for all singletons we have $m(\{x\})=0\implies\rho(\{x\})=0$ and thus $1=\mu(\{x\})=0+\lambda(\{x\})=\lambda(\{x\})$ and we can find a partition $L\cup M$ where $\lambda(M)=m(L)=0$. 
        Now, we see that for any $E$ we have $\lambda(E)=\lambda(E\cap L)=\lambda(E\cap L)+m(E\cap L)=\mu(E\cap L)\geq 0$ as $m(E\cap L)\leq m(L)=0$ proving that $\lambda$ is positive.        Then, we must have that $\lambda(M)=0\implies\lambda(C)=0$ where $C\subseteq M$ is countably infinite, we can choose such a $C$ because $M$ is required to be uncountable ($m(M)=1$.) 
        From this we see that $\infty=\mu(C)=\rho(C)+\lambda(C)=0+0$ as $m(C)=0$ forces $\rho(C)=0$ so we have a contradiction and the hypothesis was false to begin with. 
\end{enumerate}
\end{solution}
\begin{problem}\label{prob:3.14}
If $\nu$ is an arbitrary signed measure and $\mu$ is a $\sigma$-finite measure on $(X, \mM)$ such that $\nu\ll\mu$, there exists an extended $\mu$-integrable function $f:X\rightarrow[-\infty,\infty]$ such that $d\nu=f~d\mu$. Hints:
\begin{enumerate}[label=\alph*.]
	\item It suffices to assume $\mu$ is finite and $\nu$ is positive.
	\item With these assumptions there exists $E\in\mM$ that is $\sigma$-finite for $\nu$ and $\mu(E)\geq\mu(F)$ for all sets $F$ that are $\sigma$-finite for $\nu$.
	\item The radon-nikodym theorem applies on $E$. If $F\cap E=0$, then either $\nu(F)=\mu(F)=0$ or $\mu(F)>0$ and $|\nu(F)|=\infty$.
\end{enumerate}
\end{problem}
\begin{solution}
If this is true for positive $\nu$ and finite $\mu$, then for sigma finite $\mu$ we can write $X=\bigcup_1^\infty E_n$ with $E_n$ being $\mu$-finite and mutually disjoint we have $d\nu_n=f_nd\mu_n$ (we safely set $f_n=0$ outside $E_n$) where $\mu_n(E):=\mu(E\cap E_n)$ and $\nu_n(E)=\nu(E\cap E_n)$ as $\mu_n$ is finite and clearly $\nu_n\ll\mu_n$ still holds. 
Then, with $f=\sum_1^\infty f_n$ we would have $d\nu=fd\mu$. 
If $\nu$ was signed on top of that, then using \hyperref[prob:3.8]{Problem 3.8} which states $\nu\ll\mu\implies\nu^+,\nu^-\ll\mu$ we can find $f^+,f^-$ with $d\nu^-=f^-d\mu,d\nu^+=f^+d\mu$ we have $d\nu=d\nu^+-d\nu^-=(f^+-f^-)d\mu$ and $f^+-f^-$ is clearly extended integrable as $\nu$, by definition can only assume atmost one of the values $+\infty,-\infty$. 
We now prove it for this case. 
Consider the set $ \qty{\mu(H):H\in\mM,H\text{ is }\nu\,\sigma\text{-finite}}$, let $s$ be its supremum, then $s<\infty$ as $\mu$ is finite and we can find a $\nu$ $\sigma$-finite sequence of sets $F_n$ such that $\mu(F_n)\longrightarrow s$, then we can clearly see that $\mu(\cup F_n)=s$ as $\cup F_n$ is clearly still $\nu$ $\sigma$-finite and we can show both sides of the equality easily. 
Setting $E=\cup F_n$ gives us a set as described in hint (b).
The radon nikodym theorem applies on $E$ clearly. 
If $F\cap E=\emptyset$ then if we have $\nu(F)>0$ we must have $\mu(F)>0$ as $\nu\ll\mu$which forces $\nu(F)=\infty$ as otherwise the existence of $E\cup F$ contradicts the maximality of $E$ and if we have $\mu(F)>0$ then again we similarly have $\nu(F)=\infty$ such as to not contradict the maximality of $E$ so we have proved the claim in hint (c). 
There are now two cases, first one being when $\nu(E^c)=\mu(E^c)=0$, in this case we can apply radon-nikodym to get a derivative $f$ on $E$ and extend it to the rest of the set by setting it to be zero outside this, then we see that,
\[
	\nu(K)=\nu(K\cap E)+\underbrace{\nu(K\cap E^c)}_{\leq\nu(E^c)=0}=\int_{K\cap E}f\,\mathrm{d}\mu=\int_K f\,\mathrm{d}\mu
\]
as $\nu$ was assumed to be positive and $\mu(E^c)=0$ so we have $d\nu=fd\mu$ in general so we are done.
In the other case where $\mu(E^c)>0$ and $\nu(E^c)=\infty$ we can set $f$ to be $\infty$ on $E^c$, then we see that $\nu(E\cap K)=\int_{E\cap K}fd\mu$ as usual and either $\nu(E^c\cap K)=\mu(E^c\cap K)=0$ or $\nu(E^c\cap K)=\infty$ and $\mu(E^c\cap K)>0$ so we see that $\int_{E^c\cap K}fd\mu$ and $\nu(E^c\cap K)$ are either both infinite or both $0$ by part (c) of the hint as $E^c\cap K$ is disjoint from $E$, hence we have the equality,
\[
	\nu(K)=\nu(E\cap K)+\nu(E^c\cap K)=\int_{E\cap K}fd\mu+\int_{E^c\cap K}fd\mu=\int_Kfd\mu
\]
so we are done.
\end{solution}
\begin{problem}\label{prob:3.15}
A measure $\mu$ on $(X, \mM)$ is called decomposable if there is a family $\mathfrak{F}\subset\mM$ with the following properties: (i) $\mu(F)<\infty$ for all $F\in\mathfrak{F}$; (ii) the members of $\mathfrak{F}$ are disjoint and their union is $X$; (iii) if $\mu(E)<\infty$ then $\mu(E)=\sum_{F\in\mathfrak{F}}\mu(E\cap F)$; (iv) if $E\subset X$ and $E\cap F\in\mM$ for all $F\in\mathfrak{F}$ then $E\in\mM$.
\begin{enumerate}[label=\alph*.]
    \item Every $\sigma$-finite measure is decomposable.
    \item If $\mu$ is decomposable and $\nu$ is any signed measure on $(X, \mM)$ such that $\nu\ll\mu$, there exists a measurable $f:X\rightarrow[-\infty,\infty]$ such that $\nu(E)=\int_{E}f\,d\mu$ for any $E$ that is $\sigma$-finite for $\mu$, and $|f|<\infty$ on any $F\in\mathfrak{F}$ that is $\sigma$-finite for $\nu$. (Use Exercise $14$ if $\nu$ is not $\sigma$-finite.) 
\end{enumerate}
\end{problem}
\begin{solution}
\begin{enumerate}[label=\alph*.]
	\item It can be easily shown that $\mathfrak F=\{E_n:n\in\bN\}$ works where $E_n$ partition $X$ while being $\mu$-finite.
	\item Let $\mathfrak F$ be as described in the problem. 
	As $\mu(F)<\infty$ for any $F\in\mathfrak F$ we can find a measurable $f_F$ on  $F$ such that $\nu(E\cap F)=\int_{E\cap F}f_Fd\mu$ for all $E\in\mM$ by the previous problem and if $F$ is $\sigma$-finite for $\nu$ we can clearly force $|f_F|<\infty$ (in the proof presente in the book we can take all $f_j$'s which are zero outside $A_j$ to be real valued everywhere as shown in case I which makes their sum real valued too as the $A_j$ are disjoint.) as in the proof of the radon-nikodym theorem. 
	We claim that $f:=f_F$ on $F\in\mathfrak F$ is measurable and it has the properties needed. 
	For any $E\in\mathcal B_{\overline{\bR}}$ we see that $f^{-1}(E)=\bigcup_{F\in\mathfrak F}f^{-1}(E)\cap F$ and each $f^{-1}(E)\cap F=f|_{F}^{-1}(E)\cap F=f_F^{-1}(E)\cap F\in\mM$ as $f_F$ is measurable which implies that $f^{-1}(E)$ by property (iv). 
	Now, for any $E$ that is $\sigma$-finite for $\mu$, suppose $E=\cup E_n$ (disjoint union) with $\mu(E_n)<\infty$, then we see that $\mu(E_n)=\sum_{\mathfrak F}\mu(E_n\cap F)<\infty$ and it is well known that only atmost countably many $\mu(E_n\cap F)$ are nonzero, so we can find $\{F_m\}\subseteq\mathfrak F$ such that $\mu(E_n)=\sum_{m=1}^\infty\mu(E_n\cap F_m)$ and thus we also see that $\mu(E_n\setminus(\cup_m(E_n\cap F_m)))=0$ which, by absolute continuity implies $\nu(E_n)=\nu(\cup_m(F_m\cap E_n))=\sum_{m=1}^\infty\mu(E_n\cap F_m)=\sum_{m=1}^\infty\int_{E_n\cap F_m}fd\mu=\int_{E_n}fd\mu$ and thus, $\nu(E)=\sum_1^\infty\nu(E_n)=\sum_1^\infty\int_{E_n}fd\mu=\int_Efd\mu$. 
\end{enumerate}
\end{solution}
\begin{problem}\label{prob:3.16}
Suppose that $\mu, \nu$ are measures on $(X, \mM)$ with $\nu\ll\mu$, and let $\lambda=\mu+\nu$. If $f=d\nu/d\lambda$, then $0\le f<1$ $\lambda$-a.e. and $d\nu/d\mu=f/(1-f)$.
\end{problem}
\begin{solution}
(note: we treat the derivatives as equivalence classes here.) Firstly, it is obvious that $f\geq 0$ $\lambda$-a.e We have $\int_Ed\nu=\nu(E)=\int_Efd(\mu+\nu)\leq\int_Ed\nu+\int_Ed\mu=\int_E d(\mu+\nu)$, it can then be easily shown that $f\leq 1$ $\lambda$-a.e. 
If we had $f=1$ over $E$ with $\lambda(E)=\nu(E)+\mu(E)>0$ then we see that, $\nu(E)=\int_Efd\lambda=\lambda(E)$ which implies that $\mu(E)=0$ but then as $\nu\ll\mu$ we have $\nu(E)=0$ too contradicting non-nullity, this proves that $f<1$ $\lambda$-a.e. too so we have proven the first claim. 
We clearly have $\nu\ll\mu\ll\lambda$ and $\lambda\ll\mu$.
Thus, $g=d\mu/d\lambda$ exists and by Proposition $3.9$ we see that $f=d\nu/d\lambda=(d\nu/d\mu)(d\mu/d\lambda)=g(d\nu/d\mu)$ and substituting these in we get $d\nu/d\mu=f/g=f/(1-f)$ as $f+g=d\nu/d\lambda+d\mu/d\lambda=d(\mu+\nu)/d\lambda=d\lambda/d\lambda=1.$
\end{solution}
\begin{problem}\label{prob:3.17}
	Let $(X, \mM, \mu)$ be a $\sigma$-finite measure space, $\mN$ a sub-$\sigma$-algebra of $\mM$, and $\nu=\mu|_{\mN}$ is $\sigma$-finite (This condition is missing in the book statement, without which this is false: Consider the lebesgue measure on $\bR$ as the measure space and take $\mN=\{\emptyset,\bR\}, f=\chi_{[0,1]}$.) If $f\in L^{1}(\mu)$, there exists $g\in L^{1}(\nu)$ (thus $g$ is $\mN$-measurable) such that $\int_{E}f~d\mu=\int_{E}g~d\nu$ for all $E\in\mN$; if $g^{\prime}$ is another such function then $g=g^{\prime}$ $\nu$-a.e. \\
(In probability theory, $g$ is called the conditional expectation of $f$ on $\mN$.)
\end{problem}
\begin{solution}
	As $f\in L^1$ we have that $fd\mu$ is a signed measure and it is clearly still a measure restricted to $\mN$ where we see that $(fd\mu)|_{\mN}\ll\nu=\mu|_{\mN}$ and as $\mN$ is sigma finite for $\nu$ we can apply \hyperref[prob:3.14]{Problem 3.14} to deduce existence of some $g$ such that $(fd\mu)|_{\mN}=gd\nu$ which implies $\int_Efd\mu=\int_Egd\nu$ for all $E\in\mN.$ 
    Here, $g$ is extended $\nu$ integrable by construction. 
    To show that $g\in L^1(\nu)$ consider the fact that $(|f|d\mu)|_{\mN}=|(fd\mu)|_{\mN}|=|gd\nu|=|g|d\nu$ and thus, $\int_X|g|d\nu=\int_X|f|d\mu<\infty\implies g\in L^1(\nu)$. 
    The uniqueness follows from the fact that $\int_Egd\nu=\int_Eg'd\nu$ for all $E\in\mN$ implies $g=g'$ $\nu$-a.e.
\end{solution}
\begin{problem}\label{prob:3.18}
Prove Proposition 3.13c.
\end{problem}
\begin{solution}
Let $g$ be as in the definition ($d\nu=gd\mu$ with $\mu$ being positive) of $|\nu|$ then we see that $|f|d|\nu|=|f||g|d\mu$ and we also have $fd\nu=fgd\mu$ and we have that,
\[
    \qty| \int_{}^{}f\,\mathrm{d}\nu|=\qty| \int_{}^{}fg\,\mathrm{d}\mu|\leq\int|fg|d\mu=\int|f|d|\nu|
\]
\end{solution}
\begin{problem}\label{prob:3.19}
If $\nu, \lambda$ are complex measures and $\mu$ is a positive measure, then $\nu\perp\mu$ iff $|\nu|\perp|\mu|$, and $\nu\ll\lambda$ iff $|\nu|\ll\lambda$.
\end{problem}
\begin{solution}
    First, we prove $\nu\perp\mu \iff |\nu|\perp\mu$. 
    If $\nu\perp\mu$, then by definition $\nu_r\perp\mu$ and $\nu_i\perp\mu$. By \hyperref[prob:3.2]{Problem 3.2}, this implies $|\nu_r|\perp\mu$ and $|\nu_i|\perp\mu$. Using \hyperref[prob:3.9]{Problem 3.9}, their sum is also mutually singular to $\mu$, meaning $(|\nu_r|+|\nu_i|)\perp\mu$. Since total variation satisfies $|\nu| \leq |\nu_r|+|\nu_i|$, any set that is null for $|\nu_r|+|\nu_i|$ is also null for $|\nu|$. Therefore, $|\nu|\perp\mu$.
    Conversely, if $|\nu|\perp\mu$, since $|\nu_r| \leq |\nu|$ and $|\nu_i| \leq |\nu|$, we have $|\nu_r|\perp\mu$ and $|\nu_i|\perp\mu$. By \hyperref[prob:3.2]{Problem 3.2}, this implies $\nu_r\perp\mu$ and $\nu_i\perp\mu$, meaning $\nu\perp\mu$. 
    Now, we prove $\nu\ll\lambda \iff |\nu|\ll\lambda$. 
    By definition, absolute continuity with respect to a complex measure means absolute continuity with respect to its total variation, so $\nu\ll\lambda$ means $|\lambda|(E)=0 \implies \nu(E)=0$.
    If $|\nu|\ll\lambda$, then $|\lambda|(E)=0 \implies |\nu|(E)=0$. Since $|\nu(E)| \leq |\nu|(E)$, this immediately implies $\nu(E)=0$, proving $\nu\ll\lambda$.
    Conversely, suppose $\nu\ll\lambda$ and let $|\lambda|(E)=0$. Since $|\lambda|$ is a positive measure, for any measurable subset $F \subseteq E$, we also have $|\lambda|(F)=0$, which implies $\nu(F)=0$. The total variation can be equivalently defined as $|\nu|(E) = \sup \sum_{j=1}^n |\nu(E_j)|$ over all finite measurable partitions $\{E_j\}$ of $E$ by \hyperref[prob:3.21]{Problem 3.21}. Since each $E_j \subseteq E$, we have $\nu(E_j)=0$ for all $j$. Thus, the supremum is exactly $0$, giving $|\nu|(E)=0$. This proves $|\nu|\ll\lambda$.
\end{solution}
\begin{problem}\label{prob:3.20}
If $\nu$ is a complex measure on $(X,\mM)$ and $\nu(X)=|\nu|(X)$, then $\nu=|\nu|$.
\end{problem}
\begin{solution}
We see that $\nu(X)=|\nu(X)|=|\nu(E)+\nu(X\setminus E)|\leq|\nu(E)|+|\nu(X\setminus E)|\leq|\nu|(E)+|\nu|(X\setminus E)=|\nu|(X)=\nu(X)$ so all the inequalities must be equalities, which implies that $|\nu(E)|=|\nu|(E)$ as $|\nu(E)|\leq|\nu|(E)$ for all $E\in\mM.$ 
As we have $|\nu(E)|+|\nu(X\setminus E)|=|\nu(E)+\nu(X\setminus E)|$ which is real, those two complex numbers share the same arguement (unless $0$ but that case is trivial) and as $\nu(E)+\nu(X\setminus E)=\nu(X)$ is real too, they must have no complex part making them real which implies that, $\nu(E)=|\nu(E)|=|\nu|(E)$ proving the claim.
\end{solution}
\begin{problem}\label{prob:3.21}
Let $\nu$ be a complex measure on $(X, \mM)$. If $E\in\mM$ define
\[\mu_{1}(E)=\sup\qty{\sum_{1}^{n}|\nu(E_{j})|:n\in\bN,E_{1},...,E_{n} \text{ disjoint}, E=\bigcup_{1}^{n}E_{j}},\]
\[\mu_{2}(E)=\sup\qty{\sum_{1}^{\infty}|\nu(E_{j})|:E_{1},E_{2},... \text{ disjoint}, E=\bigcup_{1}^{\infty}E_{j}},\]
\[\mu_{3}(E)=\sup\qty{\qty|\int_{E}f~d\nu|:|f|\leq 1}.\]
Then $\mu_{1}=\mu_{2}=\mu_{3}=|\nu|$. (First show that $\mu_1\leq\mu_2\leq\mu_3$· To see that $\mu_3 = |\nu|$, let $f=\overline{d\nu/d|\nu|}$ and apply Proposition $3.13.$ To see that $\nu_3\leq\nu_1$, approximate $f$ by simple functions.)
\end{problem}
\begin{solution}
$\mu_1\leq\mu_2$ is obvious and $\mu_2\leq\mu_3$ holds as for any given $E$ and any partition $\qty{E_n}$ of $E$ we can define
\[
	f=\sum_n \frac{\overline{\nu(E_n)}}{|\nu(E_n)|}\cdot\chi_{E_n}
\]
where the sum is over all $n$ for which $\nu(E_n)\neq 0$, then we clearly see that $|f|\leq 1$ and $\int_Ef\,d\nu=\sum_1^\infty\int_{E_n}f\,d\nu=\sum_1^\infty|\nu(E_n)|$ as $|z|=\overline{z}z/|z|$ for nonzero $z.$ 
Clearly $\mu_3(E)\leq\int_E d|\nu|=|\nu|(E)$ by Proposition $3.13c$ and with $f$ as in the hint (this can be made to have absolute value atmost $1$ everywhere after changing it on a null set, not affecting any of the integrals below) we see that, 
\[
	\int_E\overline{ \frac{d\nu}{d|\nu|}}d\nu=\int_E\overline{ \frac{d\nu}{d|\nu|}} \frac{d\nu}{d|\nu|}d|\nu|=\int_E|f|^2d|\nu|=\int_Ed|\nu|=|\nu|(E)
\]
as $|f|=1$ $|\nu|$-a.e. so we must have $\mu_3=|\nu|$. 
Now, we can  find a sequence of simple functions $\phi_n$ converging $f$ with $|\phi_n|\leq |f|\leq 1$ and convergence can be made uniform here as $f$ is bounded on $E$, this clearly implies that $\lim|\int_E\phi_n\,d\nu|=|\int_E f\,d\nu|$ and thus, as $|\int_E\phi_n\,d\nu|\leq\mu_1(E)$ due to linearity and the triangle inequality we must have $\lim_n|\int_E f\,d\nu|\leq\mu_1(E)$ too so we are done.
\end{solution}
\begin{problem}\label{prob:3.22}
If $f\in L^{1}(\bR^{n})$, $f\ne0,$ there exist $C, R > 0$ such that $Hf(x)\ge C|x|^{-n}$ for $|x|>R$. Hence $m(\{x:Hf(x)>\alpha\})\ge C^{\prime}/\alpha$ when $\alpha$ is small, so the estimate in the maximal theorem is essentially sharp.
\end{problem}
\begin{solution}
We are given that $f$ is not $0$ a.e. and thus, $ \int_{\bR}^{}|f|\,\mathrm{d}\mu>0$ and using DCT, we see that there exists some $r>0$ such that $K=\int_{B(0,r)}^{}|f|\,\mathrm{d}\mu>0$. 
Then, for all $x$, as $B(0,r)\subseteq B(x,|x|+r)$ we have
\[
Hf(x)\geq \frac{1}{m(B(x,|x|+r))}\int_{B(x,|x|+r)}|f|\,d\mu\geq \frac{1}{c\cdot(|x|+r)^n}\int_{B(0,r)}|f|\,d\mu\geq \frac{K}{(|x|+r)^n}
\]
where $c=m(B(0,1))$. 
Thus, we can find some $C>0$ and $R>0$ like in the statement. 
Hence, for small $\alpha>0$ we have $\qty{x:Hf(x)>\alpha}\supseteq \qty{x:C\alpha^{-1}\geq|x|^n>R^n}$. 
From this we can deduce that, $m( \qty{x:Hf(x)>\alpha})\geq c\cdot(C\alpha^{-1}-R^n)$ which is atleast $C'\alpha^{-1}$ for some $C'>0$ and small $\alpha.$
\end{solution}
\begin{problem}\label{prob:3.23}
A useful variant of the Hardy-Littlewood maximal function is $H^{*}f(x)=\sup\{\frac{1}{m(B)}\int_{B}|f(y)|dy:B \text{ is a ball and } x\in B\}$. Show that $Hf\le H^{*}f\le2^{n}Hf$.
\end{problem}
\begin{solution}
If $x\in B(p,r)=B$ then we must have that $B(p,r)\subseteq B(x,|x-p|+r)$, thus
\[
	\frac{1}{m(B)}\int_B|f|\,dm\leq \frac{c\cdot(|x-p|+r)^n}{c\cdot r^n}Hf(x)\leq 2^nHf(x)
\]
where $c=m(B(0,1))$, as $|x-p|<r$ so we must have $H^*f\leq 2^nHf$ and $H^*f\geq Hf$ is trivially true. 
\end{solution}
\begin{problem}\label{prob:3.24}
If $f\in L_{loc}^{1}$ and $f$ is continuous at $x$, then $x$ is in the Lebesgue set of $f$.
\end{problem}
\begin{solution}
	For all $\epsilon>0$ we can find $\delta>0$ such that $f(B(x,\delta)) \subseteq B(f(x),\epsilon)$, thus for all $0<\delta'<\delta$ we must have $m(B(x,\delta'))^{-1}\int_{B(x,\delta')}|f(x)-f(y)|dy\leq\epsilon$ so it must converge to $0$ implying $x\in L_f.$
\end{solution}
\begin{problem}\label{prob:3.25}
If $E$ is a Borel set in $\bR^{n}$, the density $D_{E}(x)$ of $E$ at $x$ is defined as $D_{E}(x)=\lim_{r\rightarrow0}\frac{m(E\cap B(r,x))}{m(B(r,x))}$ whenever the limit exists.
\begin{enumerate}[label=\alph*.]
    \item Show that $D_{E}(x)=1$ for a.e. $x\in E$ and $D_{E}(x)=0$ for a.e. $x\in E^{c}$.
    \item Find examples of $E$ and $x$ such that $D_{E}(x)$ is a given number $\alpha\in(0,1)$, or such that $D_{E}(x)$ does not exist.
\end{enumerate}
\end{problem}
\begin{solution}
\begin{enumerate}[label=\alph*.]
\item Consider the borel measure $\nu(F):=\mu(E\cap F)$, then we have a lebesgue decomposition $\nu=0+\chi_Edm$ and Theorem $3.22.$ we have that $D_E(x)=\chi_E(x)$ for $m$-a.e. $x.$ from which the conclusion follows.
\item Consider $E=\bigcup_{k=1}^{\infty}(B(0,2^{-2k+1})\setminus B(0,2^{-2k}))$ and $x=0$. Let $c = m(B(0,1))$. Then we see that for $m\in\bN$ we have 
    \[ m(E\cap B(0,2^{-2m+1})) = c\sum_{k=m}^{\infty}\left(2^{n(-2k+1)}-2^{-2kn}\right) = c(2^n-1)\frac{2^{-2mn}}{1-2^{-2n}}. \]
    Dividing by $m(B(0, 2^{-2m+1})) = c 2^{n(-2m+1)}$, the density along this sequence of radii is:
    \[ \frac{m(E\cap B(0,2^{-2m+1}))}{m(B(0,2^{-2m+1}))} = \frac{c(2^n-1)2^{-2mn}/(1-2^{-2n})}{c 2^{-2mn} 2^n} = \frac{2^n}{2^n+1}. \]
    Now consider the sequence of radii $r = 2^{-2m}$. The ball $B(0, 2^{-2m})$ intersects $E$ starting from index $k=m+1$, so:
    \[ m(E\cap B(0,2^{-2m})) = c\sum_{k=m+1}^{\infty}\left(2^{n(-2k+1)}-2^{-2kn}\right) = c(2^n-1)\frac{2^{-2n(m+1)}}{1-2^{-2n}}. \]
    Dividing by $m(B(0, 2^{-2m})) = c 2^{-2mn}$ yields:
    \[ \frac{m(E\cap B(0,2^{-2m}))}{m(B(0,2^{-2m}))} = \frac{c(2^n-1)2^{-2mn}2^{-2n}/(1-2^{-2n})}{c 2^{-2mn}} = \frac{1}{2^n+1}. \]
    Since $\frac{2^n}{2^n+1}\neq\frac{1}{2^n+1}$ for all $n \geq 1$, two subsequences tend to two different limits, hence $D_E(0)$ does not exist. 
    To find an example where $D_E(x) = \alpha \in (0,1)$, choose $x=0$. 
	If $n \geq 2$, define $E$ as a sector of $\bR^n$ originating at $0$. Let $A \subset S^{n-1}$ be a measurable set on the unit sphere such that its surface measure is $\alpha \sigma(S^{n-1})$ which we can do by considering the function $f:[-1,1]\to\bR$ defined by $f(t):=\sigma( \qty{x\in S^{n-1}:x_1\leq t})$ which can be showed to be continuous via lower and upper continuity, after which we can use the intermediate value theorem. Define $E = \{ y \in \bR^n \setminus \{0\} : y/|y| \in A \}$. For any $r > 0$, $m(E \cap B(0,r)) = \alpha m(B(0,r))$, meaning $D_E(0) = \alpha$ identically for all $r>0$.
    If $n = 1$, define $E = \bigcup_{k=1}^\infty \left( \left[ \frac{1}{k} - \frac{\alpha}{k(k+1)}, \frac{1}{k} \right] \cup \left[ -\frac{1}{k}, -\frac{1}{k} + \frac{\alpha}{k(k+1)} \right] \right)$. For $r = 1/k$, $m(E \cap (-r, r)) = 2 \sum_{j=k}^\infty \frac{\alpha}{j(j+1)} = \frac{2\alpha}{k} = \alpha (2r)$. Bounding the measure for $r \in (\frac{1}{k+1}, \frac{1}{k})$ shows the ratio to $\alpha$, hence $D_E(0) = \alpha$. 
\end{enumerate}
\end{solution}
\begin{problem}\label{prob:3.26}
If $\lambda$ and $\mu$ are positive, mutually singular Borel measures on $\bR^{n}$ and $\lambda+\mu$ is regular, then so are $\lambda$ and $\mu$.
\end{problem}
\begin{solution}
Define $\nu=\lambda+\mu$ and let $\lambda(A)=\mu(A^c)=0$. 
For compact $K$ we see that $\lambda(K)+\mu(K)=\nu(K)<\infty$, which implies that $\lambda(K),\mu(K)<\infty$. 
First, we show outer regularity. 
If $\lambda(E) = \infty$, the condition $\lambda(E) = \inf\{\lambda(U) : U \supseteq E, U \text{ open}\}$ is trivially satisfied by any open $U \supseteq E$. 
Assume $\lambda(E) < \infty$. 
We wish to find an open $U \supseteq E$ with $\lambda(U \setminus E) < \epsilon$. 
We split $E$ into $E \cap A^c$ and $E \cap A$. 
Notice that $\nu(E \cap A^c) = \lambda(E \cap A^c) + \mu(E \cap A^c) = \lambda(E \cap A^c) + 0 \le \lambda(E) < \infty$. 
Since $\nu(E \cap A^c) < \infty$ and $\nu$ is regular, there exists an open $U_1 \supseteq (E \cap A^c)$ with $\nu(U_1 \setminus (E \cap A^c)) < \epsilon/2$, which implies $\lambda(U_1 \setminus (E \cap A^c))\leq\nu(U_1\setminus(E\cap A^c))< \epsilon/2$. 
For the remaining piece $E \cap A$, we have $\lambda(E \cap A) \le \lambda(A) = 0$. 
We can write $A = \bigcup_{k=1}^\infty A_k$ where $A_k = A \cap B(0,k)$. 
As each $A_k$ is bounded, $\nu(A_k) < \infty$, so there exist open sets $V_k \supseteq A_k$ with $\nu(V_k \setminus A_k) < \epsilon/2^{k+1}$. 
This implies $\lambda(V_k) = \lambda(V_k \setminus A_k) + \lambda(A_k) < \epsilon/2^{k+1} + 0$. 
Taking $U_2 = \bigcup_{k=1}^\infty V_k$ gives an open set containing $A$ (and thus $E \cap A$) with $\lambda(U_2) \le \sum \lambda(V_k) < \epsilon/2$. 
Let $U = U_1 \cup U_2$. 
Then $U$ is open, $U \supseteq E$, and $U \setminus E \subseteq (U_1 \setminus (E \cap A^c)) \cup U_2$. 
Thus, $\lambda(U \setminus E) \le \lambda(U_1 \setminus (E \cap A^c)) + \lambda(U_2) < \epsilon/2 + \epsilon/2 = \epsilon$ so we have the second condition. 
The first condition is trivial.
\end{solution}
\begin{problem}\label{prob:3.27}
Verify the assertions in Examples 3.25.
\end{problem}
\begin{solution}
This is trivial.
\end{solution}
\begin{problem}\label{prob:3.28}
If $F\in NBV$, let $G(x)=|\mu_{F}|((-\infty,x])$. Prove that  $|\mu_{F}|=\mu_{T_{F}}$ by showing that $G=T_{F}$ via the following steps.
\begin{enumerate}[label=\alph*.]
	\item From the definition of $T_{F}$, $T_{F}\le G$.
	\item $|\mu_{F}(E)|\le\mu_{T_{F}}(E)$ when $E$ is an interval, and hence when $E$ is a Borel set.
    \item $|\mu_{F}|\le\mu_{T_{F}}$, and hence $G\le T_{F}$.
\end{enumerate}
\end{problem}
\begin{solution}
\begin{enumerate}[label=\alph*.]
    \item For any $x\in\bR$ and $x_0<\ldots<x_n=x$ we see that,
        \[
            \sum_{k=1}^n|F(x_k)-F(x_{k-1})|=\sum_{k=1}^n|\mu_F((x_{k-1},x_k])|\leq\sum_{k=1}^n|\mu_F|((x_{k-1},x_k])=|\mu_F|((x_0,x])
        \]
       Which is clearly less than $G(x).$ 
       This shows that $T_F\leq G$ as $T_F$ is the supremum over all sums of this form.
   \item For open (we only need these) intervals, the inequality follows from the fact that $T_F(x)-T_F(y)\geq|F(x)-F(y)|$ for $x\geq y$ as we can take limits and say that, for $a<b$
       \[
           \mu_{T_F}((a,b))=T_F(b-)-T_F(a)\geq|F(b-)-F(a)|=|\mu_F((a,b))|
       \]
       (keep in mind that all the one sided limits exist in this case.) 
       We know by construction that $\mu_F=\mu_1-\mu_2+i(\mu_3-\mu_4)$ for finite lebesgue stjeltes measures $\mu_i$ as in chapter $1.$ 
       Let $\mu_5=\mu_{T_F},$ another lebesgue stjeltes measure. 
       Given borel $E$, by  regularity we can find open sets $U_n^{(i)}\supseteq E$ such that $\mu_i(U_n^{(i)})\to\mu_i(E)$ for all $i$. 
       Notice that by the \twemoji{sandwich} theorem and monotonicty of $\mu_i$, the sequence of open sets $V_n:=\cap_i U_n^{(i)}$ have the same convergence property by monotonicity of measures and the sandwich property. 
       As $\mu_i$ are all finite we see that $\mu_F(V_n)\to\mu_F(E).$ 
       If $\cup I_n$ is a countable disjoint union of open intervals we see that $|\mu_F(\cup I_n)|\leq\mu_{T_F}(\cup I_n)$ by the triangle inequality and sigma additivity of measures.
       From this, we can conclude that $|\mu_F(V_n)|\leq\mu_{T_F}(V_n)$ for all $n$ and we have the desired inequality by taking the limit as $n\to\infty$
   \item This can be proven by  the previous part and \hyperref[prob:3.21]{Problem 3.21} directly (look at $\mu_1$) in the hyperlinked problem. 
       From this we thus have that $G\leq T_F$ by considering left open rays.
\end{enumerate}
Now that we have $G=T_F$ we must have that $\mu_G=\mu_{T_F}$ and $\mu_G$ is none other than $|\mu_F|.$
\end{solution}
\begin{problem}\label{prob:3.29}
If $F\in NBV$ is real-valued, then $\mu_{F}^{+}=\mu_{P}$ and $\mu_{F}^{-}=\mu_{N}$ where $P$ and $N$ are the positive and negative variations of $F$.
\end{problem}
\begin{solution}
It is easily seen that, the positive and negative variations are nonnegative for $NBV$ functions as $F(-\infty)=T_F(-\infty)=0.$ 
Now, we clearly have $\mu_F=\mu_{P}-\mu_{N}$ (this is true for h-intervals and the sets where this holds forms a sigma algebra due to these being finite measures) which proves that $\mu_P\geq\mu_F^+$ and $\mu_N\geq\mu_F^-$ by \hyperref[prob:3.4]{Problem 3.4}. 
Now, from page 103 we can also see that
\[
    P(x)=\sup\qty{\sum_{k=1}^n[F(x_k)-F(x_{k-1})]^+:x_0<\ldots<x_n=x} + 0
\]
as $F\in NBV$, now notice that for $a<b$, we have 
\[
    [F(b)-F(a)]^+=[\mu_F((a,b])]^+=[\mu_F^+((a,b])-\mu_F^-((a,b])]^+\leq[\mu_F^+((a,b])]^+=\mu_F^+((a,b])
\]
From this we can easily see that $\mu_P((-\infty,x])=P(x)\leq\mu_F^+((-\infty,x])$. 
But earlier we proved that $\mu_P\geq\mu_F^+$ so we must have equality above and hence $\mu_P=\mu_F^+$ can be easily shown to hold. 
\end{solution}
\begin{problem}\label{prob:3.30}
Construct an increasing function on $\bR$ whose set of discontinuities is $\bQ$.
\end{problem}
\begin{solution}
    Consider an ennumeration $\{q_n\}$ of the rationals. 
    Define $f(x):=\sum_{q_n<x}2^{-n}$ then it is obvious that for rationals $r$ we have $f(q_n+)-f(q_n-)\geq 2^{-n}$ making it discontinuous on the rationals. 
    Now, if $x$ is an irrational, let $\delta>0$ be such that $q_1,\ldots,q_m\not\in(x-\delta,x+\delta)$, we can do this becasue $x$ is irrational. 
    In this case, we see that for $y\in(x-\delta,x+\delta)$
    \[
        |f(x)-f(y)|=\sum_{q_n\in[x-\delta,x+\delta)}\frac1{2^n}\leq\sum_{n>m}\frac1{2^n}=\frac1{2^{m+1}}
    \]
    Thus we have continuity at the irrationals as this can be made arbitrarily small for larger $m$ and small enough $\delta>0.$
\end{solution}
\begin{problem}\label{prob:3.31}
Let $F(x)=x^{2}\sin(x^{-1})$ and $G(x)=x^{2}\sin(x^{-2})$ for $x\ne0,$ and $F(0)=G(0)=0$.
\begin{enumerate}[label=\alph*.]
    \item $F$ and $G$ are differentiable everywhere (including $x=0$).
    \item $F\in BV([-1,1])$, but $G\notin BV([-1,1])$.
\end{enumerate}
\end{problem}
\begin{solution}
    The first part is an exercise in calculus and the next follows from $F'$ being bounded ( $|F'|\leq3$ to be a bit more specific) and $T_G(1)-T_G(n^{-1/2})\gg\log n$ (this is not to be interpreted as $\infty-\infty$ but rather the total variation in the smaller interval $(n^{-1/2},1) )$
\end{solution}
\begin{problem}\label{prob:3.32}
If $F_{1},F_{2},...,F\in NBV$ and $F_{j}\rightarrow F$ pointwise, then $T_{F}\le \liminf T_{F_{j}}$.
\end{problem}
\begin{solution}
For any $x\in\bR$ and $\epsilon>0$ consider some $x_0<\ldots<x_n=x$ such that,
\[
    \sum_{k=1}^n|F(x_k)-F(x_{k-1})|\geq T_F(x)-\epsilon
\]
Then, by pointwise convergence we see that for large enough $j$,
\[
    -\epsilon+\sum_{k=1}^n|F(x_k)-F(x_{k-1})|\leq\sum_{k=1}^n|F_j(x_k)-F_j(x_{k-1})|\leq T_{F_j}(x)
\]
This shows that $T_{F}(x)-2\epsilon\leq T_{F_j}(x)$ for large enough $j$ and thus, $T_F(x)-2\epsilon\leq\liminf T_{F_j}(x)$. 
As $\epsilon>0$ was arbitrary we are done.
\end{solution}
\begin{problem}\label{prob:3.33}
If $F$ is increasing on $\bR$, then $F(b)-F(a)\ge\int_{a}^{b}F^{\prime}(t)dt$.
\end{problem}
\begin{solution}
As we only care about things being defined a.e. we define $F_n(x):=n(F(x+1/n)-F(x))\geq 0$ and we see that, $F_n\to F'$ a.e. thus, by fatou's lemma we have
\[
    \int_a^bF'\leq\liminf\int_a^b F_n=\liminf_nn\int_a^bF(x+1/n)-F(x)\,\mathrm{d}m(x)
\]
If we set $G(x):=\int_a^xF(t)dm(t)$ then we can pass to rieman integrablity (monotone functions are rieman integrable) and see that the last expression above is
\[
    \liminf_n\frac{G(b+1/n)-G(a+1/n)-(G(b)-G(a))}{1/n}=G'(b)-G'(a)=F(b)-F(a)
\]
As $G'(t)=F(t)$ by the fundamental theorem of calculus, when $F$ is continuous at $a,b.$ 
Now, as $F$ is continuous a.e. the set of points where it is continuous must be dense and we can thus find sequences $a_n\searrow a<b\nwarrow b_n$ where $F$ is continuous and thus, 
\[
    F(b)-F(a)\geq F(b-)-F(a+)=\lim_{n\to\infty} F(b_n)-F(a_n)\geq\lim\int_{[a_n,b_n]}F'=\int_{(a,b)}F'=\int_a^bF'
\]
where the last two equalities are due to the fact that $A\mapsto\int_AF'$ is a measure here ($(a_n,b_n)$ increases to $(a,b)$) and $m(\{a,b\})=0$.
\end{solution}
\begin{problem}\label{prob:3.34}
Suppose $F, G\in NBV$ and $-\infty<a<b<\infty.$
\begin{enumerate}[label=\alph*.]
    \item By adapting the proof of Theorem 3.36, show that
    \[\int_{[a,b]}\frac{F(x)+F(x-)}{2}dG(x)+\int_{[a,b]}\frac{G(x)+G(x-)}{2}dF(x)=F(b)G(b)-F(a-)G(a-)\]
    \item If there are no points in $[a, b]$ where $F$ and $G$ are both discontinuous, then
    \[\int_{[a,b]}F~dG+\int_{[a,b]}G~dF=F(b)G(b)-F(a-)G(a-).\]
\end{enumerate}
\end{problem}
\begin{solution}
\begin{enumerate}[label=\alph*.]
    \item Following steps similar to those the proof of theorem 3.36, taking $\Omega=\{(x,y):a\leq x\leq y\leq b\}$ we can see that,
        \[
            \int_{[a,b]}FdG-F(a-)[G(b)-G(a-)]=G(b)[F(b)-F(a-)]-\int_{[a,b]}G(x-)dF(x)
        \]
        We can also switch $F,G$ and get another equation
        \[
            \int_{[a,b]}GdF-G(a-)[F(b)-F(a-)]=F(b)[G(b)-G(a-)]-\int_{[a,b]}F(x-)dG(x)
        \]
        rearranging and adding these we get,
        \[
            \int_{[a,b]}(F(x)+F(x-))dG(x)+\int_{[a,b]}(G(x)+G(x-))dF(x)=2(F(b)G(b)-F(a-)G(a-))
        \]
    \item If $F$ is continuous at $x$ then we see that 
        \[
            \mu_F(\{x\})=\lim_n\mu_F((x+1/n,x])=\lim_n F(x)-F(x+1/n)=0
        \]
        and similarly for $G$. 
        As $F,G$ are NBV we can write them as the differnece of two increasing functions which only ever have countably many discontinuities and thus the set of discontinuities here form countable sets which we call $D_F,D_G$ for $F$ and $G$ respectively. 
        Then we can see that $\mu_F(D_G)=\sum_{x\in D_G}\mu_F(\{x\})=\sum_{x\in D_G}0=0$ because $F$ is continuous on $D_G$ and similarly $\mu_G(D_F)=0$. 
        Now, in the equation from the previous part we can write the LHS as, 
        \[
            \int_{[a,b]\setminus D_F}(F(x)+F(x-))dF(x)+\int_{[a,b]\setminus D_G}(G(x)+G(x-))dF(x)
        \]
        as the integrands are continuous over their domain of integration, we see that $F(x)+F(x-)=2F(x)$ and similarly for $G.$ 
        From this we can conclude that, 
        \[
            \int_{[a,b]}2FdG+\int_{[a,b]}2GdF=2(F(b)G(b)-F(a-)G(a-))
        \]
        from which we get the desired equality after cancelling the $2$ out.
\end{enumerate}
\end{solution}
\begin{problem}\label{prob:3.35}
If $F$ and $G$ are absolutely continuous on $[a, b]$, then so is $FG$, and
\[\int_{a}^{b}(FG^{\prime}+GF^{\prime})(x)dx=F(b)G(b)-F(a)G(a).\]
\end{problem}
\begin{solution}
It suffices to show that $F$ being absolutely continuous implies $F^2$ is absolutely continuous. 
Indeed, if $F$ is absolutely continuous then it is continuous too and thus is bounded on the compact set $[a,b]$, say $|F|\leq M$. 
If, $\epsilon>0$ is arbitrary, we can find $\delta>0$ with,
\[
    \sum_{i=1}^n(b_i-a_i)<\delta\implies\sum_{i=1}^n|F(b_i)-F(a_i)|<\frac{\epsilon}{M+1}\text{ for }a\leq a_1<b_1<\ldots<b_n\leq b
\]
Then, 
\[
    \sum_{i=1}^n|F(b_i)^2-F(a_i)^2|\leq M\sum_{i=1}^n|F(b_i)-F(a_i)|\leq\frac{\epsilon M}{M+1}<\epsilon
\]
so $F^2$ is also absolutely continuous. 
Now, it is clear that $F,G$ being absolutely continuous makes their product absolutely continuous too. 
From Theorem 3.35c we see that $(FG)'\in L^1([a,b],m)$ and the same also holds for $F',G'$. 
As they all exist a.e., they exist together a.e. as well and we can say that $(FG)'=FG'+F'G$ a.e. and again by the cited theorem,
\[
    F(b)G(b)-F(a)G(a)=\int_a^b(FG)'(x)dx=\int_a^b(FG'+F'G)(x)dx
\]
\end{solution}
\begin{problem}\label{prob:3.36}
Let $G$ be a continuous increasing (we need $G$ to be strictly increasing, otherwise consider constant functions.) function on $[a, b]$ and let $G(a)=c$, $G(b)=d.$
\begin{enumerate}[label=\alph*.]
    \item If $E\subset[c,d]$ is a Borel set, then $m(E)=\mu_{G}(G^{-1}(E)).$ (First consider the case where $E$ is an interval.)
    \item If $f$ is a Borel measurable and integrable function on $[c, d]$, then $\int_{c}^{d}f(y)dy=\int_{a}^{b}f(G(x))dG(x)$. In particular, $\int_{c}^{d}f(y)dy=\int_{a}^{b}f(G(x))G^{\prime}(x)dx$ if $G$ is absolutely continuous.
    \item The validity of (b) may fail if $G$ is merely right continuous rather than continuous.
\end{enumerate}
\end{problem}
\begin{solution}
\begin{enumerate}[label=\alph*.]
    \item The case where this is an interval follows as shown, for $c\leq x<y\leq d$ we just have $\mu_G(G^{-1}([x,y]))=\mu_G([G^{-1}(x),G^{-1}(y)])=G\circ G^{-1}(y)-G\circ G^{-1}(x)=y-x=m([x,y])$. 
        We can easily show that the collection of sets $E$ for which this works is a sigma algebra as the inverse commutes with the complement and unions, we thus have that the borel sigma algebra over $[c,d]$ is contained in this one which suffices.
    \item If $f=\chi_E$ ($E\in\mathcal B_{[c,d]}$) then, $\int_a^b\chi_E(G(x))dG(x)=\int_a^b\chi_{G^{-1}(E)}dG(x)=\mu_G(G^{-1}(E))=m(E)=\int_c^d\chi_Edm$ and we can conclude by linearity of the integral and MCT as usual. 
       If $G$ is NBV, we have that $dG=G'dm$ (can be proven from 3.35c) and by \hyperref[prob:2.14]{Problem 2.14} we are done (split the integrable function into its two parts as usual.)
   \item In that case the range of $f$ needn't be an interval in the first place. 
       For example, take $f(x):=x\chi_{[0,1/2)}(x)+(x+1)\chi_{[1/2,1]}$ which is right continuous, increasing and maps $[0,1]$ to $[0,1/2)\cup[3/2,2].$
\end{enumerate}
\end{solution}
\begin{problem}\label{prob:3.37}
Suppose $F:\bR\rightarrow\bC$. There is a constant $M$ such that $|F(x)-F(y)|\le M|x-y|$ for all $x,y\in\bR$ (that is, $F$ is Lipschitz continuous) iff $F$ is absolutely continuous and $|F^{\prime}|\le M$ a.e.
\end{problem}
\begin{solution}
    For the only if direction, absolute continuity follows directly from the definition by considering $\epsilon/{M+1}$ instead of $\epsilon$ and for the derivative to exist a.e. (it is trivially bounded as needed when it does exist) we consider intervals $(n,n+1]$ for $n\in\bZ.$ 
    Clearly $F$ is still absolutely continuous on this small interval and by Theorem 3.35c we again see that $F'$ exists a.e. in $(n,n+1]$ and thus it exists a.e. in $\bR$ too as the union of countably many disjoint null sets is still null.
    To prove the if part, for $x<y$ consider $F$ on the interval $[x,y].$ 
    By Theorem 3.35c we can write $F(x)-F(y)=\int_x^yF'(t)dt$ again and by hypothesis we have,
    \[
        |F(x)-F(y)|=\left|\int_x^yF'(t)dt\right|\leq\int_x^y|F'(t)|dt\leq M(y-x)=M|x-y|
    \]
    The case of $x>y$ is no different.
\end{solution}
\begin{problem}\label{prob:3.38}
If $f:[a,b]\rightarrow\bR$, consider the graph of $f$ as a subset of $\bC$, namely, $\{t+if(t): t\in[a,b]\}$. The length $L$ of this graph is by definition the supremum of the lengths of all inscribed polygons.
\begin{enumerate}[label=\alph*.]
    \item Let $F(t)=t+if(t);$ then $L$ is the total variation of $F$ on $[a, b]$.
    \item If $f$ is absolutely continuous, $L=\int_{a}^{b}[1+f^{\prime}(t)^{2}]^{1/2}dt$.
\end{enumerate}
\end{problem}
\begin{solution}
\begin{enumerate}[label=\alph*.]
    \item This is trivial, just remember how $|x+iy|=\sqrt{x^2+y^2}.$
    \item If $f$ is absolutely continuous then so is $F$ and by problem 3.35c it can be shown that $\mu_F= F'dm$ on this interval and thus, $\mu_{T_F}=|F'|dm$ by \hyperref[prob:3.28]{Problem 3.28}, this shows that
        \[
            L=T_F(b)-T_F(a)=\mu_{T_F}((a,b])=\int_{(a,b]}|F'(t)|dm(t)=\int_a^b[1+f'(t)^2]^{1/2}dm(t)
        \]
        as $F'(t)=1+i\cdot f'(t)$ (when it exists.)
\end{enumerate}
\end{solution}
\begin{problem}\label{prob:3.39}
    If $\{F_{j}\}$ is a sequence of nonnegative increasing functions on $[a, b]$ such that $F(x)=\sum_{1}^{\infty}F_{j}(x)<\infty$ for all $x\in[a,b]$, then $F^{\prime}(x)=\sum_{1}^{\infty}F_{j}^{\prime}(x)$ for a.e. $x\in[a,b]$. (It suffices to assume $F_j\in NBV$. onsider the meausres $\mu_{F_j}.)$
\end{problem}
\begin{solution}
    Keep in mind the following lemma : If $P_1,P_2,\ldots$ hold a.e. then $\land P_i$ does too, the proof is left as an exercsie.
    We can consider the functions $G_j(x)=F_j(x+)-F_j(a)$ and $G(x)=F(x+)-F(a)$ instead. These have the same properties as $F_j$ and $F$, and we extend them to make them $NBV$ as done before (they were trivially in $BV([a,b])$ already). Note that $G=\sum G_j$ everywhere, which follows from an application of the monotone convergence theorem with respect to the counting measure. 
    Consider the measures $\mu_G$ and $\mu_{G_j}$, for which it holds that $\mu_G=\sum_j\mu_{G_j}$ since they agree on right-closed rays. Now, consider the Lebesgue-Radon-Nikodym decompositions $dG = g \, dm + d\lambda$ and $dG_j = g_j \, dm + d\lambda_j$. 
    We have $\lambda_j \perp m$, and thus $\sum \lambda_j \perp m$ too by \hyperref[prob:3.9]{Problem 3.9} (notice that $\lambda_j,\lambda$ are positive measures as well.)
    Since $\mu_G = \sum \mu_{G_j}$, we must have $dG = (\sum_j g_j) \, dm + \sum \lambda_j$. By the uniqueness criterion in the Lebesgue-Radon-Nikodym theorem, this implies $g = \sum g_j$ a.e. and $\lambda = \sum \lambda_j$. 
    It is clear from the lebesgue differentiation theorem that $g_j = G_j'$ a.e. and $g = G'$ a.e., so the lemma at the beginning implies that $G' = \sum G_j'$ a.e.  
    If we now show that $F' = G'$ and $F_j' = G_j'$ a.e., then we would be done using the same lemmat. Suppose $p$ and $q$ are such that $p = q$ a.e. and both are differentiable a.e. We want to show that $p' = q'$ a.e. We can find a set $A \subseteq [a,b]$ where $p = q$ and both are differentiable, such that $m([a,b] \setminus A) = 0$ by the lemma. 
    Notice that $A$ must be dense in this interval, so for any $x \in A$, we can find a sequence $x_n \in A \setminus \{x\}$ that tends to $x$. By the sequential criterion for limits, we have
    \[
        p'(x) = \lim_{n \to \infty} \frac{p(x_n) - p(x)}{x_n - x} = \lim_{n \to \infty} \frac{q(x_n) - q(x)}{x_n - x} = q'(x)
    \]
    for all $x \in A$, so we are done.
\end{solution}
\begin{problem}\label{prob:3.40}
Let $F$ denote the Cantor function on $[0, 1]$ (see §1.5), and set $F(x)=0$ for $x<0$ and $F(x)=1$ for $x>1$. Let $\{[a_{n},b_{n}]\}$ be an enumeration of the closed subintervals of $[0, 1]$ with rational endpoints, and let $F_{n}(x)=F((x-a_{n})/(b_{n}-a_{n}))$. Then $G=\sum_{1}^{\infty}2^{-n}F_{n}$ is continuous and strictly increasing on $[0, 1]$, and $G^{\prime}=0$ a.e. (Use exercise 39)
\end{problem}
\begin{solution}
$F_n$ are cleary increasing, continuous and nonnegative and we also see that $G=\sum F_n\leq1<\infty$ as $F_n\leq 1$ for all $n$ by definition. 
$G$ is also continuous by weierstrass' M-test being the limit of an uniformly converging sequence of continuous functions. 
Given any $0\leq x<y\leq 1$ we can find some interval with rational endpoints $[a_n,b_n]\subset(x,y)$ and we see that $G(y)-G(x)\geq 2^{-n}(F_n(y)-F_n(x))=2^{-n}(1-0)>0$ by definition again, making $G$ strictly increasing. 
We see that $F$ is constant on intervals outside the cantor set and thus has derivative $0$ here and as the cantor set has lebesgue measure $0$ we can show that $F'=0$ a.e. by standard procedures. 
Now, by \hyperref[prob:3.39]{Problem 3.39} we see that $G'(x)=\sum 2^{-n}F_n'(x)$ a.e. and as $F_n'=0$ a.e. too (can be easily shown), we must have $G'=0$ a.e.
\end{solution}
\begin{problem}\label{prob:3.41}
Let $A\subset[0,1]$ be a Borel set such that $0<m(A\cap I)<m(I)$ for every subinterval $I$ of $[0, 1]$.
\begin{enumerate}[label=\alph*.]
    \item Let $F(x)=m([0,x]\cap A)$. Then $F$ is absolutely continuous and strictly increasing on $[0, 1]$, but $F^{\prime}=0$ on a set of positive measure.
    \item Let $G(x)=m([0,x]\cap A)-m([0,x]\setminus A).$ Then $G$ is absolutely continuous on $[0, 1]$, but $G$ is not monotone on any subinterval of $[0, 1]$.
\end{enumerate}
\end{problem}
\begin{solution}
\begin{enumerate}[label=\alph*.]
    \item Clearly $F$ is lipschitz continuous with constant $1$ and thus its absolutely continuous by \hyperref[prob:3.37]{Problem 3.37} and its strictly increasing as $F(y)-F(x)=m((x,y]\cap A)>0$ for any $x<y$ by definition. 
        As $F$ is absolutely continuous, by theorem 3.35c we see that $F'\geq 0$ exists a.e, is integrable and $F(x)=F(x)-F(0)=\int_0^xF'(t)dm(t)$. 
       $F$ can be extended to $\bR$ by setting it equal to $0=F(0)$ for $x<0$ and $m(A)=F(1)$ for $x>1$ and this is NBV so we can talk about $\mu_F.$ 
       We also see that $\mu_F((-\infty,x]))=(F'dm)((-\infty,x])$ for all $x$ and thus, $\mu_F=F'dm$. 
       If $E$ is a set of positive meausre and we have $\mu_F(E)=\int_EF'dm=0$ then it must be true that $F'=0$ a.e. in $E$ which again furnishes a set of positve meausre (same as $E$) on which $F'=0.$ 
       Notice that the measure defined by $\nu(E)=m(E\cap A)$ is the same as $\mu_F$ as they agree on right closed rays and hence we just need to find some $E$ such that $m(E)>0$ whereas $\nu(E)=0$ and the perfect (not a topological property!) candidate is $A^c$ for which $m(A^c)=1-m(A)>0$ whereas $\nu(A^c)=m(A^c\cap A)=0$ so we are done.
   \item We can easily see that $m(A^c\cap I)=m(I)-m(A\cap I)\in(0,m(I))$ as well, if we let $H(x)=m([0,x]\cap A^c)$ then we see that $H$ has the same properties as $F$ from the previous part. 
       Thus, $G$ is also absolutely continuous being the linear combination of two absolutely continuous functions. 
       Given any interval $[x,y]\subset[0,1]$ ($x<y$) suppose for the sake of contradiction that WLOG. $G$ is increasing in this interval, this implies that $G'\geq 0$ a.e. here but this can not be true as it implies $\mu_G=\mu_F-\mu_H$ has $[x,y]$ as a positive set. 
      But we clearly see that for $V=A^c\cap [x,y]\subseteq[x,y]$ we have $\mu_F(V)=m(V\cap A)=0$ whereas $\mu_H(V)=m(V\cap A^c)>0$ and thus, $\mu_G(V)<0$ which is a contradiction and thus we are done.
\end{enumerate}
\end{solution}
\begin{problem}\label{prob:3.42}
A function $F:(a,b)\rightarrow\bR(-\infty\le a<b\le\infty)$ is called convex if $F(\lambda s+(1-\lambda)t)\le\lambda F(s)+(1-\lambda)F(t)$ for all $s, t\in(a,b)$ and $\lambda\in(0,1)$.
\begin{enumerate}[label=\alph*.]
    \item $F$ is convex iff for all $s, t, s', t' \in (a, b)$ such that $s\le s^{\prime}<t^{\prime}$ and $s<t\le t^{\prime}$,
    \[\frac{F(t)-F(s)}{t-s}\le\frac{F(t^{\prime})-F(s^{\prime})}{t^{\prime}-s^{\prime}}\]
    \item $F$ is convex iff $F$ is absolutely continuous on every compact subinterval of $(a, b)$ and $F^{\prime}$ is increasing (on the set where it is defined).
    \item If $F$ is convex and $t_{0}\in(a,b)$, there exists $\beta\in\bR$ such that $F(t)-F(t_{0})\ge \beta(t-t_{0})$ for all $t\in(a,b)$.
    \item (Jensen's Inequality) If $(X, \mM, \mu)$ is a measure space with $\mu(X)=1$, $g:X\rightarrow(a,b)$ is in $L^{1}(\mu)$, and $F$ is convex on $(a, b)$, then
    \[F\left(\int g~d\mu\right)\le\int F\circ g~d\mu.\]
    (let $t_0=\int gd\mu$ and $t=g(x)$ in (c), and integrate.)
\end{enumerate}
\end{problem}
\begin{solution}
\begin{enumerate}[label=\alph*.]
    \item This is a standard result from introductory real analysis, it can be proved by just manipulating the inequalities.
    \item If $F$ is convex then we see that on any compact subinterval $[c,d]$, for $x,y\in[c,d]$ if $F'$ exists on $x,y$ ($x<y$) then, for $h>0$ small enough, by part (a)
        \[
            \frac{F(x+h)-F(x)}{(x+h)-x}\leq\frac{F(v)-F(u)}{v-u}\leq\frac{F(y+h)-F(y)}{(y+h)-y}
        \]
        where $x<x+h<y<y+h$ and as $h\to 0$ from above, as the derivatives exist at $x,y$ we get that $F'(x)\geq F'(y).$ 
        It is another standard result that convex functions are continuous and we can also see that if, $c=s'<t'$ and $c<t\leq t'$ then,
        \[
            \frac{F(t)-F(c)}{t-c}\leq \frac{F(t')-F(c)}{t'-c}
        \]
        Thus, $G(x):=(F(x)-F(c))/(x-c)$ is increasing, similarly we can show that $H(x)=(F(d)-F(x))/(d-x)$ is decreasing. 
        This, implies that $G(c+),H(d-)$ both exist and we can show that these bound all slopes using an arguement similar to the one we used to show that $F'$ is increasing on the set its defined. 
        We thus have, for all $x,y$ such that $c\leq x<y\leq d$ we have
        \[
            G(c+)\leq\frac{F(y)-F(x)}{y-x}\leq H(d-)
        \]
        and thus $F$ is lipschitz with constant atmost $\max\{|G(c+)|,|H(d-)|\}$ so it is absolutely continuous by a previous problem. 
        We now prove the if direction. 
        By hypothesis, for any such $s,t,s',t'$ as in part (a), we work in the compact subinterval $[s,t]$ where the inequality in part (a) reduces to the following by theorem 3.35c,
        \[
            \frac1{t-s}\int_s^tF'(t)dt\leq\frac1{t'-s'}\int_{s'}^{t'}F'(t)dt
        \]
        where $s\leq s'<t'$ and $s<t\leq t'.$ 
        This is trivial if $s'>t$ so we assume otherwise that $s'\leq t$, where this inequality is true as, 
        \[
            \frac1{t-s}\int_s^tF'= \frac{s'-s}{t-s}\cdot\underbrace{\frac1{s'-s}\int_s^{s'}F'}_{\leq\frac1{t-s'}\int_{s'}^tF'}+ \frac{t-s'}{t-s}\cdot\frac1{t-s'}\int_{s'}^tF'\leq\frac1{t-s'}\int_{s'}^tF'
        \]
        The inequality in the underbrace is true as all the values appearing in the integralin the left are atmost those appearing in the right and we average these. 
        It can be similarly shown that the expression on the right is atmost $\frac1{t'-s'}\int_{s'}^{t'}F'$ so we are done. 
    \item Just like we did for $c,d$ in part (b), it can be showed that the left and right derivatives exist at every point and we pick any $\beta\in[F'(t_0-),F'(t_0+)]$ (left and right derivatives and not the one sided limits of our derivative.) 
        The conclusion follows just as we did in part (b), we show part of it. 
        \[
            F'(t_0+)=\lim_{h\searrow 0}\frac{F(t_0+h)-F(t_0)}{h}\leq\frac{F(y)-F(t_0)}{y-t_0}
        \]
        for any $y>t_0$ as $t_0+h<y$ eventually. 
    \item We see that $\int_Xgd\mu\in(a,b)$ so from part (c), we can find some constant $c$ such that,
        \[
            F(t)-F\qty(\int gd\mu)\geq c\cdot\qty(t-\int gd\mu)
        \]
        for all $t\in(a,b)$ and thus we can substitute $t=g(x)$ to get,
        \[
            F\circ g(x)-F\qty(\int gd\mu)\geq c\cdot\qty(g(x)-\int gd\mu)
        \]
        for all $x\in X.$ 
        Now, integratint both sides we see that,
        \[
            \int F\circ g\,d\mu-F\qty(\int g\,d\mu)\geq c\cdot\int gd\mu-\int gd\mu=0
        \]
        so we are done.
\end{enumerate}
\end{solution}
\end{document}
